% New formalization incorporated into revision III; originals preserved.
% Source owners: U016, equations extractor-diferencias-k and lector-armonico-*;
% registro_k.tex, subsections k-pi-h, k-hadamard, and k-transversal.
\subsection{Integral compatibility and determination of the transverse register}
\label{img09:seccion}

The domain of this construction is the representation of the register of an
APP--TRIT--TPK history. We distinguish evaluation of the register's events,
compatibility of its channels, and inversion of its
coordinates. The following result gives a complete integral criterion for
the second level and an explicit factorization connecting the first to it.
The proof uses the register operators and establishes their
relations for every element of the stated integer domains.

\subsubsection{Characterization by adjacent increments}

Indices belong to $\mathbb Z/12\mathbb Z$. Let
\[
 (Sz)_i=z_{i+1},\qquad D_3=I-S^3,\qquad D_4=I-S^4,
 \qquad Q(z)=\sum_{i=0}^{11}z_i.
\]
We retain the transverse extractor
\[
 \mathcal D:\mathbb Z^{12}\longrightarrow\mathbb Z^{25},
 \qquad k\longmapsto(D_3k,D_4k,Q(k)).
\]
For a triple of integer coordinates $(b,c,q)$, write
\begin{equation}
 w_i=c_i-b_{i+1},\qquad
 P_0=0,\quad P_i=\sum_{j=0}^{i-1}w_j\ (1\leq i\leq12),
 \qquad T(w)=\sum_{i=0}^{11}P_i.
 \label{img09:incremento}
\end{equation}

\begin{theorem}[Integral image of the transverse extractor]
\label{img09:imagen}
A triple $(b,c,q)\in\mathbb Z^{25}$ belongs to
$\operatorname{im}\mathcal D$ if and only if it satisfies
\begin{align}
 b_i&=w_i+w_{i+1}+w_{i+2}&& (0\leq i<12),
 \label{img09:relaciones}\\
 \sum_{i=0}^{11}w_i&=0,
 \label{img09:cierre}\\
 q+T(w)&\equiv0\pmod {12}.
 \label{img09:congruencia}
\end{align}
Its unique integer preimage is
\begin{equation}
 k_i=\frac{q+T(w)}{12}-P_i,\qquad0\leq i<12.
 \label{img09:inversa}
\end{equation}
The twelve equations in \eqref{img09:relaciones} and equation
\eqref{img09:cierre} are linearly independent over $\mathbb Q$.
\end{theorem}
\begin{proof}
For an output of $\mathcal D$, we have
\[
 c_i-b_{i+1}=(k_i-k_{i+4})-(k_{i+1}-k_{i+4})=k_i-k_{i+1}.
\]
The sum of three increments gives $b_i$; the cyclic sum of all of them
is zero. Moreover, $P_i=k_0-k_i$, so
$q=12k_0-T(w)$. This yields all the conditions and the inverse formula.

Conversely, \eqref{img09:congruencia} makes formula
\eqref{img09:inversa} integer-valued; \eqref{img09:cierre} allows its
cyclic extension. Its adjacent differences are $w_i$. Therefore,
$D_3k=b$ and
\[
 (D_4k)_i=w_i+w_{i+1}+w_{i+2}+w_{i+3}
 =w_i+b_{i+1}=c_i.
\]
Summing the inverse gives $Q(k)=q$. The same formula proves uniqueness.
For independence, the invertible integer change of variables
$(b,c)\mapsto(b,w=c-Sb)$ converts the first twelve relations into
$b=(I+S+S^2)w$. Each solves for a different coordinate of $b$.
The condition $\sum w_i=0$ is an additional independent relation.
\end{proof}

The rational image has dimension twelve. Its integral structure adds
a congruence of order twelve: this is the constructive manifestation of the
last Smith factor of $\mathcal D$. The eleven increments
$w_0,\ldots,w_{10}$ and the charge $q$ are sufficient coordinates, with
$w_{11}=-\sum_{i=0}^{10}w_i$ and congruence
\eqref{img09:congruencia}. The original twenty-five coordinates
retain these same twelve degrees of freedom through thirteen relations.

\subsubsection{Commutativity with the harmonic readout}

Let $\mathcal H_{12}$ be the Hadamard transformation of the register,
with orbits $(r,r+3,r+6,r+9)$ for $r=0,1,2$. Its block is
\[
 H_4=\begin{pmatrix}
 1&1&1&1\\1&1&-1&-1\\1&-1&1&-1\\1&-1&-1&1
 \end{pmatrix},\qquad H_4^2=4I_4.
\]
For compatible $(b,c,q)$, define
\[
 B_r=\sum_{j=0}^3c_{r+3j},\qquad
 A_0=\frac{q+2B_0+B_1}{3},\quad A_1=A_0-B_0,
 \quad A_2=A_1-B_1.
\]
The reader $\Pi_H$ in \eqref{eq:k-bloques-firmados} has blocks
\[
 (\Pi_H(b,c,q))^{(r)}=
 (A_r,b_{r+3}-b_{r+9},b_r+b_{r+6},b_r-b_{r+6}).
\]

\begin{proposition}[Agreement of the two reconstructions]
\label{img09:triangulo}
On $\operatorname{im}\mathcal D$, we have
\[
 \Pi_H\mathcal D=\mathcal H_{12},\qquad
 \mathcal D^{-1}=\tfrac14\mathcal H_{12}\Pi_H.
\]
The image of $\Pi_H$ restricted to that domain is exactly
\[
 \mathcal L_H=
 \{u\in\mathbb Z^{12}:\mathcal H_{12}u\equiv0\pmod4\}.
\]
\end{proposition}
\begin{proof}
For $k=\mathcal D^{-1}(b,c,q)$, let
$a_r=\sum_{j=0}^3k_{r+3j}$. The shift by four positions
takes each orbit to the next, so $B_r=a_r-a_{r+1}$.
Together with $q=a_0+a_1+a_2$, these identities give $A_r=a_r$.
On one orbit, write $(x_0,x_1,x_2,x_3)$ for the coordinates
of $k$ and $d_j=x_j-x_{j+1}$. Then
\[
 d_1-d_3=x_0+x_1-x_2-x_3,\quad
 d_0+d_2=x_0-x_1+x_2-x_3,\quad
 d_0-d_2=x_0-x_1-x_2+x_3.
\]
These are the three remaining rows of $H_4k^{(r)}$. Inversion and the
integral characterization follow from $\mathcal H_{12}^2=4I$.
\end{proof}

The condition $\Pi_H(b,c,q)\in\mathcal L_H$ controls only the
harmonic output. Membership of the twenty-five coordinates in
$\operatorname{im}\mathcal D$ additionally requires
\eqref{img09:relaciones}--\eqref{img09:congruencia}.
For example, adding to $c$ the vector $e_0-e_3$ preserves the three orbital
sums and leaves $\Pi_H$ unchanged. Starting from $(0,0,0)$ thus
gives a triple with $\Pi_H=0$ that violates
\eqref{img09:relaciones}. This is an exact negative control for
channel compatibility, independent of any terminal evaluation.

\subsubsection{Determination and transport on register orbits}

Theorem~\ref{img09:imagen} determines $k$ once its channels
and charge have been produced. As a condition on outputs not yet
individualized, the same image is a group isomorphic to
$\mathbb Z^{12}$. With only a charge $q$ fixed, the solutions
form an affine coset of
\[
 L_0=\{v\in\mathbb Z^{12}:\sum v_i=0\}
 =\bigoplus_{i=0}^{10}\mathbb Z(e_i-e_{11}).
\]
For example, $100\mathbf1+e_0$ and $100\mathbf1+e_1$ are two
distinct registers of charge $1201$, with components in
$\{0,\ldots,999\}$. Each satisfies all image conditions
and all Hadamard congruences through its own
channels. These witnesses compare the compatibility equations;
they are not assigned the status of terminal HMT histories.

\begin{proposition}[Cyclic covariance and the register stabilizer]
\label{img09:covariancia}
Let $\rho$ be the window translation on a domain of registers
$\mathscr R$, and let $S$ be the preceding translation on $\mathbb Z^{12}$.
On the orbit of $R\in\mathscr R$, a covariant reader is
determined by a vector
\[
 k\in\operatorname{Fix}(S^p),
\]
where $p\mid12$ is the least period of $R$ under $\rho$.
The charge condition $Q(k)=q$ admits integer solutions if and only if
$12/p$ divides $q$; when they exist, they form an affine coset of rank
$p-1$. In particular, on a free orbit, covariance and charge
leave eleven degrees of freedom.
\end{proposition}
\begin{proof}
The rule $F(\rho^jR)=S^jk$ is well defined exactly when
$S^pk=k$. Such vectors repeat $12/p$ times a word of
length $p$. Their charge is $\frac{12}{p}\sum_{i=0}^{p-1}k_i$.
Divisibility and rank follow from this formula. Channel
covariance follows from $D_aS=SD_a$ and $QS=Q$.
On $\mathcal L_H$, the corresponding action is
$\widehat S_H=\mathcal H_{12}S\mathcal H_{12}/4$;
by Proposition~\ref{img09:triangulo}, both reconstructions
intertwine these actions.
\end{proof}

If the domain carries a marked origin, one must use the action that
also transports that mark. The proposition concerns this declared
cyclic action; it does not presuppose periodicity of the enriched
history. Naturality with respect to truncations beyond the cut at
$108$ is obtained, for any already defined reader $F_{108}$, through
$F_N=F_{108}\circ\tau_{N,108}$. This naturality preserves the initial
evaluation rule; it does not by itself choose among the possible rules
on the first one hundred and eight events.

\subsubsection{Eventwise evaluation and composition of contributions}

The preceding characterization provides a factorization of extraction through two
effective readouts of the register:
\begin{equation}
 \omega:\mathscr R\longrightarrow
 \{w\in\mathbb Z^{12}:\sum w_i=0\},\qquad
 q_{\rm reg}:\mathscr R\longrightarrow\mathbb Z,
 \quad q_{\rm reg}(R)+T(\omega(R))\equiv0\pmod {12}.
 \label{img09:contrato-productor}
\end{equation}
These maps must be evaluated on the events, orientation,
incidence, and memory produced by APP--TRIT--TPK. Once their
actions on those data are given, the composition
\[
 R\longmapsto(w,q)\longmapsto
 ((I+S+S^2)w,(I+S+S^2+S^3)w,q)
 \xrightarrow{\Pi_H}u\xrightarrow{\mathcal H_{12}/4}k
\]
is determined and satisfies all inverse checks. The formula
$w=c-Sb$ also connects a producer that already supplies both
channels: it is a check on its output, preceding any comparison
with the terminal witness.

\begin{proposition}[Accumulation of contributions and prospective uniqueness]
\label{img09:acumulacion}
Suppose that the initial state and each event $a_t$ of the register
produce, through rules fixed beforehand, contributions
$(\delta w_t,\delta q_t)$ with
\[
 \sum_i\delta w_{t,i}=0,\qquad
 \delta q_t+T(\delta w_t)\equiv0\pmod {12}.
\]
With compatible initial conditions $(w^{(0)},q^{(0)})$, the
update
\[
 w^{(t+1)}=w^{(t)}+\delta w_t,\qquad
 q^{(t+1)}=q^{(t)}+\delta q_t
\]
produces a unique register at each stage. Its increment is
\[
 \delta k_{t,i}=
 \frac{\delta q_t+T(\delta w_t)}{12}-P_i(\delta w_t).
\]
The construction respects concatenation of registers, and the readout
of a prefix remains unchanged under subsequent prolongations.
\end{proposition}
\begin{proof}
The compatibility conditions are preserved under addition. Formula
\eqref{img09:inversa} is additive in $(w,q)$ and produces the
stated increment. Induction fixes each state successively.
The sum over two segments is the concatenated sum, and a prefix receives
exactly the same summands before and after the history is prolonged.
\end{proof}

The proposition provides a sufficient criterion, not a requirement of
linearity for every TPK reader. If events are transported by nontrivial
actions, their contributions can be expressed in the terminal chart
through the corresponding affine cocycle. The cocycle equation then
preserves the same independence from the partition of the trajectory.
In both cases, the event rules and the initial state constitute the
constructive data that must come from the article: choosing them to
reproduce an expected vector would replace that provenance with a different task.

The result of this development is the exact determination condition
\eqref{img09:contrato-productor}, its proven inverse, and its transport
identities. Its scope is the transverse register; the other
HMT constructions retain their own domains and proofs.
